\Needspace{8\baselineskip}
\section{Comparison with the supplied geometry papers}
\label{sec:geometry}

\subsection{What the analogy explains}
The supplied paper \emph{Living on the edge} studies phase transitions for
randomly oriented convex cones and convex recovery problems \citep{almt2014}.
For a closed convex cone $C\subset\R^d$, its statistical dimension is
\[
 \delta(C)=\mathbb E\norm{\Pi_C(g)}^2,
 \qquad g\sim N(0,I_d),
\]
which extends linear dimension. If $C,K$ are closed convex cones and $Q$ is
a Haar-uniform orthogonal transformation, Theorem I gives, for
$0<\eta<1$ and $a_\eta=\sqrt{8\log(4/\eta)}$,
\begin{align*}
 \delta(C)+\delta(K)\le d-a_\eta\sqrt d
 &\ \Longrightarrow\ \mathbb P(C\cap QK\ne\{0\})\le\eta,\\
 \delta(C)+\delta(K)\ge d+a_\eta\sqrt d
 &\ \Longrightarrow\ \mathbb P(C\cap QK\ne\{0\})\ge1-\eta.
\end{align*}
This identifies both a transition location and a quantitative window through
an intrinsic geometric parameter. The theorem is not merely an empirical
fit to a transition curve.

The supplied \emph{Sharp Phase Transitions in Euclidean Integral Geometry}
develops concentration for weighted intrinsic-volume distributions of convex
bodies and applies it to integral-geometric questions
\citep[Theorems I--II and Section 4]{lotz2026}. Its central rotation and
rigid-motion volumes play the role of location parameters. This suggests
a productive conceptual comparison: identify a meaningful random variable,
prove concentration, and then link its typical value to feasibility.

Random $k$-SAT has a discrete version of a random covering problem. A clause
forbids precisely $2^{n-k}$ vertices of the Boolean cube; UNSAT means the
union of the forbidden subcubes covers every vertex. This perspective appears
directly in Carenini's proof. But its subcubes are sparse, coordinate-aligned
and strongly overlap. There is no Haar-random rotation, no given convex cone,
and no established statistical-dimension formula for $\alpha_k$ in this
comparison.

\subsection{Why the obvious convex relaxation loses the transition}
A clause with positive-variable set $P$ and negative-variable set $N$ can
be written exactly on Boolean points as
\[
 \sum_{i\in P}x_i+\sum_{j\in N}(1-x_j)\ge1.
\]
It is tempting to replace the Boolean cube by its convex hull and apply a
convex feasibility theorem. The following elementary observation shows why
that first step cannot work.

\begin{proposition}[The elementary clause relaxation cannot detect UNSAT]
\label{prop:lp}
Let $F$ be a CNF formula in which every clause contains at least two
literals. Replace each positive literal by $x_i$, each negative literal by
$1-x_i$, impose that each clause's literal sum is at least $1$, and relax
$x\in\{0,1\}^n$ to $x\in[0,1]^n$. This linear system is feasible for every
such $F$. In particular, feasibility of this relaxation is not equivalent
to satisfiability, even for proper $k$-CNF with fixed $k\ge2$.
\end{proposition}
\begin{proof}
Set every $x_i=1/2$. Every literal then has value $1/2$, so a clause with
$r\ge2$ literals has sum $r/2\ge1$. For an explicit UNSAT proper $k$-CNF,
take all $2^k$ sign patterns on the same $k$ variables. Every Boolean
assignment falsifies the clause with the opposite signs, so no assignment
satisfies the formula, although the same fractional point is feasible.
\end{proof}


The obstruction is specific to this relaxation. It does not rule out stronger
lifts, semidefinite or moment constraints, or a geometric description that
retains integrality. Such a program would first need a precise object, a
proved link between its geometric event and Boolean satisfiability, and the
correct distribution of orientations or constraints. Concentration of the
wrong feasible set would not resolve the random-$k$-SAT problem.

\noindent\textbf{Provenance.} Proposition~\ref{prop:lp} is an elementary
expository observation, recorded as P001 with state \texttt{proof-written};
it is not presented as novel or independently reviewed. The geometric
comparison supplies a direction for further study, not a proved transfer
of either supplied geometric theorem.
